\chapter{Stress Tensor, Weyl Structure, and Improvement}
\label{ch:volume-iii-stress-tensor-weyl}

\section{The Stress Tensor from Background Metrics}

Let a Euclidean local theory be coupled to a smooth background metric
\(g_{\mu\nu}\) on a \(D\)-dimensional manifold \(M\).  The stress tensor is
defined by the first variation of the action or of the renormalized generating
functional with respect to \(g_{\mu\nu}\).  In an action formulation,
\begin{equation}
  \delta S[\Phi;g]
  =
  \frac12
  \int_M \dd^D x\,\sqrt g\,
  T^{\mu\nu}(x)\,\delta g_{\mu\nu}(x),
  \label{eq:stress-tensor-metric-variation}
\end{equation}
where \(\Phi\) denotes the dynamical fields and
\(g=\det(g_{\mu\nu})\).  With this sign convention, a flat-space translation
current is obtained by evaluating \(T_{\mu\nu}\) at
\(g_{\mu\nu}=\delta_{\mu\nu}\).

If the theory is invariant under diffeomorphisms of the background, then an
infinitesimal vector field \(\xi^\mu\) gives
\[
  \delta_\xi g_{\mu\nu}
  =
  \nabla_\mu\xi_\nu+\nabla_\nu\xi_\mu .
\]
Substituting this variation into
\eqref{eq:stress-tensor-metric-variation} and integrating by parts gives
\begin{equation}
  \nabla_\mu T^{\mu\nu}=0
  \label{eq:stress-tensor-covariant-conservation}
\end{equation}
as an operator equation away from insertions.  In flat space this becomes
\(\partial_\mu T^{\mu\nu}=0\).  Symmetry of \(T_{\mu\nu}\) follows from
variation with respect to the symmetric tensor \(g_{\mu\nu}\).

\section{Weyl Transformations}

A Weyl transformation is a local rescaling of the background metric:
\[
  g_{\mu\nu}(x)\mapsto \ee^{2\omega(x)}g_{\mu\nu}(x),
\]
where \(\omega\) is a smooth real function.  Infinitesimally,
\[
  \delta_\omega g_{\mu\nu}=2\omega g_{\mu\nu}.
\]
The corresponding metric variation gives
\[
  \delta_\omega S
  =
  \int_M\dd^D x\,\sqrt g\,
  \omega(x)\,T^\mu{}_\mu(x).
\]
Thus Weyl invariance of the flat-space theory coupled to background metrics
is expressed by vanishing trace, modulo contact terms and possible curvature
anomalies:
\begin{equation}
  T^\mu{}_\mu=0
  \label{eq:weyl-trace-zero}
\end{equation}
in flat space.  The use of a background metric is not an additional dynamical
field; it is a device for defining stress-tensor insertions and their Ward
identities.

In even dimensions, the renormalized generating functional may transform
under a Weyl rescaling by a local curvature functional.  Such a Weyl anomaly
does not obstruct flat-space conformal Ward identities at separated points,
but it contributes contact terms and curved-space response data.  The core
monograph records this structural fact; the special two-dimensional extension
of conformal symmetry is deferred to a separate treatment.

\section{Improvement of the Stress Tensor}

The stress tensor is not unique.  If \(L_{\mu\nu}(x)\) is a local symmetric
operator, then
\begin{equation}
  T_{\mu\nu}
  \longmapsto
  T'_{\mu\nu}
  =
  T_{\mu\nu}
  +
  \left(
  \partial_\mu\partial_\nu-\delta_{\mu\nu}\partial^2
  \right)L
  \label{eq:scalar-stress-improvement}
\end{equation}
for a scalar \(L\) preserves conservation in flat space.  More general
improvements are possible, but \eqref{eq:scalar-stress-improvement} is the
case needed for the free scalar and many scalar fixed points.  Its trace is
\[
  T'^\mu{}_\mu
  =
  T^\mu{}_\mu
  -(D-1)\partial^2L .
 \]
If \(T^\mu{}_\mu\) is a total Laplacian of a local scalar, an improvement can
make the trace vanish in flat space.

The improvement can be represented by adding a curvature coupling to the
background-metric action.  For a scalar field \(\phi\), the term
\[
  \Delta S_\xi
  =
  \frac{\xi}{2}
  \int_M\dd^D x\,\sqrt g\,R(g)\phi^2
\]
vanishes in flat space when \(R(g)=0\), but its metric variation contributes
to \(T_{\mu\nu}\) at flat \(g_{\mu\nu}\).  This is why a term invisible in the
flat-space action can change the flat-space stress tensor.

\section{The Conformally Coupled Free Scalar}

Consider a real massless scalar field in Euclidean signature with action
\[
  S_0[\phi]
  =
  \frac12\int\dd^D x\,
  \partial_\mu\phi\,\partial^\mu\phi .
\]
The canonical stress tensor is
\[
  T_{\mu\nu}^{(0)}
  =
  \partial_\mu\phi\,\partial_\nu\phi
  -
  \frac12\delta_{\mu\nu}(\partial\phi)^2 .
\]
Its trace is
\[
  T^{(0)\mu}{}_\mu
  =
  \left(1-\frac D2\right)(\partial\phi)^2 .
\]
Using the equation of motion \(\partial^2\phi=0\), the identity
\[
  \partial^2\phi^2
  =
  2(\partial\phi)^2+2\phi\,\partial^2\phi
\]
shows that the trace is a Laplacian modulo the equation of motion.

The improved stress tensor is
\begin{equation}
  T_{\mu\nu}
  =
  \partial_\mu\phi\,\partial_\nu\phi
  -
  \frac12\delta_{\mu\nu}(\partial\phi)^2
  +
  \xi_\ast
  \left(
  \delta_{\mu\nu}\partial^2-\partial_\mu\partial_\nu
  \right)\phi^2,
  \label{eq:improved-free-scalar-stress}
\end{equation}
where
\begin{equation}
  \xi_\ast=\frac{D-2}{4(D-1)}.
  \label{eq:conformal-scalar-xi}
\end{equation}
Indeed,
\[
  T^\mu{}_\mu
  =
  \left(1-\frac D2+2(D-1)\xi_\ast\right)(\partial\phi)^2
  +
  2(D-1)\xi_\ast\,\phi\,\partial^2\phi,
\]
and this vanishes as an operator equation modulo the scalar equation of
motion.  The corresponding curved-space action is
\begin{equation}
  S[\phi;g]
  =
  \frac12
  \int_M\dd^D x\,\sqrt g\,
  \left[
  g^{\mu\nu}\partial_\mu\phi\partial_\nu\phi
  +
  \xi_\ast R(g)\phi^2
  \right].
  \label{eq:conformal-scalar-curved-action}
\end{equation}

\begin{figure}[H]
  \centering
  \resizebox{0.90\textwidth}{!}{%
  \begin{tikzpicture}[x=1cm,y=1cm,every node/.style={font=\scriptsize}]
    \node[font=\small] at (0,2.2) {flat action};
    \draw[rounded corners=4pt,thick] (-1.8,0.45) rectangle (1.8,1.45);
    \node at (0,0.95) {\(\frac12\int(\partial\phi)^2\)};

    \draw[-{Latex[length=2mm]},thick] (2.05,0.95)--(3.45,0.95);
    \node[align=center] at (2.75,1.45) {couple to\\metric};

    \node[font=\small] at (5.45,2.2) {curvature term};
    \draw[rounded corners=4pt,thick] (3.65,0.45) rectangle (7.25,1.45);
    \node at (5.45,0.95) {\(\frac{\xi_\ast}{2}\int\sqrt g\,R\phi^2\)};

    \draw[-{Latex[length=2mm]},thick] (7.50,0.95)--(8.90,0.95);
    \node[align=center] at (8.20,1.45) {vary\\\(g_{\mu\nu}\)};

    \node[font=\small] at (10.65,2.2) {improved tensor};
    \draw[rounded corners=4pt,thick] (9.05,0.45) rectangle (12.25,1.45);
    \node at (10.65,0.95) {\(T^\mu{}_\mu=0\)};
    \node[blue!65!black,align=center] at (5.45,-0.45)
      {the curvature term affects flat-space\\stress-tensor insertions};
  \end{tikzpicture}
  }
  \caption{A curvature coupling can vanish on the flat background and still
  change the flat-space stress tensor, because the stress tensor is obtained by
  varying the background metric before setting it flat.}
  \label{fig:stress-improvement-curvature}
\end{figure}

\section{Ward Identities with Stress-Tensor Insertions}

Let \(\mathcal O_i(x_i)\) be local operators inserted at separated points.
The Ward identity for diffeomorphisms has the local form
\[
  \partial_\mu
  \left\langle
  T^{\mu\nu}(x)\mathcal O_1(x_1)\cdots\mathcal O_n(x_n)
  \right\rangle
  =
  -\sum_{i=1}^n
  \delta(x-x_i)
  \partial_{x_i^\nu}
  \left\langle
  \mathcal O_1(x_1)\cdots\mathcal O_n(x_n)
  \right\rangle
  +
  \cdots ,
\]
where the omitted contact terms account for spin indices and for the
variation of composite operators under diffeomorphisms.  The trace Ward
identity similarly records the response under Weyl transformations.  At a
flat-space conformal fixed point, the separated-point trace insertion
vanishes, while contact terms encode scaling dimensions and conformal
transformations of the inserted operators.

The charge form of the Ward identity follows by integrating
\(T^\mu{}_\nu\epsilon^\nu\) over a small sphere around each insertion.  This
is the bridge between the local stress-tensor equation and the conformal
algebra acting on local operators.
